\section{Phạm vi nghiên cứu}
\label{sec:scope}

Luận văn tập trung vào phạm vi nghiên cứu cụ thể sau đây, với các giới hạn và giả định rõ ràng:

\subsection{Phạm vi về bài toán}

\textbf{Phân loại nhị phân lối sống:} Luận văn tập trung vào bài toán phân loại nhị phân (binary classification) để phân biệt thực khuẩn thể độc lực (virulent/lytic) và thực khuẩn thể ôn hòa (temperate/lysogenic). Các lối sống phức tạp hơn hoặc các trạng thái trung gian không được xem xét trong phạm vi nghiên cứu này.

\textbf{Dữ liệu đầu vào:} Phương pháp được thiết kế để hoạt động trên các contig DNA được lắp ráp từ dữ liệu metagenomic, không phải trên các đoạn đọc thô (raw reads). Các contig được giả định là đã được xác định là có nguồn gốc từ phage thông qua các công cụ nhận dạng phage như VirSorter, VIBRANT, hoặc các phương pháp tương tự.

\subsection{Phạm vi về dữ liệu}

\textbf{Phạm vi độ dài contig:} Nghiên cứu tập trung vào bốn nhóm độ dài contig:
\begin{itemize}
    \item Nhóm A: 100-400 bp (contig rất ngắn)
    \item Nhóm B: 400-800 bp (contig ngắn)
    \item Nhóm C: 800-1200 bp (contig trung bình)
    \item Nhóm D: 1200-1800 bp (contig dài)
\end{itemize}

Các contig ngắn hơn 100 bp hoặc dài hơn 1800 bp không được đánh giá chi tiết trong nghiên cứu này, mặc dù phương pháp về mặt kỹ thuật có thể xử lý các chuỗi có độ dài tùy ý nhờ ALiBi positional encoding của DNABERT-2.

\textbf{Nguồn dữ liệu:} Tập dữ liệu được xây dựng từ hai nghiên cứu gần đây: DeePhage \cite{wu2021deephage} và DeepPL \cite{zhang2024deeppl}, bao gồm 2,241 hệ gen phage hoàn chỉnh (707 temperate, 1,534 virulent). Các contig được tạo ra bằng phương pháp sliding window để mô phỏng kịch bản metagenomic thực tế.

\textbf{Giới hạn về đa dạng sinh học:} Mặc dù tập dữ liệu bao gồm nhiều họ phage khác nhau, nó vẫn chỉ đại diện cho một phần nhỏ của sự đa dạng phage trong tự nhiên. Hiệu suất của mô hình trên các phage thuộc các họ hiếm gặp hoặc chưa được nghiên cứu có thể khác với kết quả báo cáo.

\subsection{Phạm vi về phương pháp}

\textbf{Mô hình nền tảng:} Nghiên cứu sử dụng DNABERT-2 \cite{zhou2023dnabert} làm mô hình nền tảng genome. Các mô hình nền tảng khác như Nucleotide Transformer \cite{dalla2025nucleotide}, HyenaDNA, hoặc các biến thể BERT khác không được đánh giá trong phạm vi luận văn này.

\textbf{Kiến trúc downstream:} Luận văn tập trung vào kiến trúc lai kết hợp CNN đa tỷ lệ và attention pooling. Các kiến trúc downstream khác như RNN, LSTM, hoặc pure transformer decoder không được khảo sát chi tiết.

\textbf{Phương pháp so sánh:} Luận văn so sánh với bốn phương pháp tiên tiến: DNABERT-2, PhaTYP, DeePhage, và ProkBERT. Các phương pháp học máy truyền thống như PHACTS, PhageAI, và BACPHLIP không được đánh giá lại vì đã được chứng minh có hiệu suất thấp hơn đáng kể trên dữ liệu contig trong các nghiên cứu trước.

\subsection{Phạm vi về đánh giá}

\textbf{Metrics đánh giá:} Luận văn sử dụng ba metrics chính: sensitivity (recall), specificity, và accuracy. Các metrics khác như precision, F1-score, AUC-ROC, hoặc Matthews correlation coefficient không được báo cáo chi tiết để duy trì tính nhất quán với các nghiên cứu trước đó.

\textbf{Phương pháp validation:} Sử dụng stratified 5-fold cross-validation để đánh giá hiệu suất. Không thực hiện đánh giá trên các tập dữ liệu độc lập bên ngoài (external validation) do hạn chế về tài nguyên và tính sẵn có của dữ liệu có nhãn chất lượng cao.

\textbf{Phân tích ablation:} Luận văn thực hiện phân tích so sánh giữa PhaBERT-CNN và DNABERT-2 baseline để đánh giá tác động của các thành phần CNN và attention. Các nghiên cứu ablation chi tiết hơn (ví dụ: loại bỏ từng nhánh CNN, thay đổi kernel sizes, so sánh các cơ chế attention khác nhau) không được thực hiện trong phạm vi luận văn này.

\subsection{Phạm vi về triển khai}

\textbf{Framework và thư viện:} Mô hình được triển khai sử dụng PyTorch và thư viện Transformers của Hugging Face. Các framework khác như TensorFlow/Keras không được xem xét.

\textbf{Tài nguyên tính toán:} Các thí nghiệm được thực hiện trên GPU NVIDIA RTX A4000 16GB. Phân tích về khả năng mở rộng (scalability) trên các cấu hình phần cứng khác hoặc tối ưu hóa cho triển khai trên CPU không được đề cập chi tiết.

\textbf{Công cụ phần mềm:} Luận văn không phát triển một công cụ phần mềm độc lập (standalone software tool) hoặc web server. Mã nguồn được cung cấp dưới dạng script Python để tái tạo kết quả nghiên cứu.

\subsection{Giới hạn và giả định}

\textbf{Giả định về chất lượng dữ liệu:} Các contig đầu vào được giả định là:
\begin{itemize}
    \item Đã được xác định chính xác là có nguồn gốc từ phage
    \item Không chứa lỗi giải trình tự đáng kể
    \item Không bị ô nhiễm bởi DNA vi khuẩn chủ
\end{itemize}

\textbf{Giới hạn về tổng quát hóa:} Hiệu suất của mô hình được đánh giá trên dữ liệu mô phỏng (simulated data) được tạo từ hệ gen hoàn chỉnh. Hiệu suất trên dữ liệu metagenomic thực tế từ các môi trường đa dạng có thể khác biệt và cần được xác thực thêm.

\textbf{Không xem xét các yếu tố sinh học phức tạp:} Luận văn không xem xét các trường hợp đặc biệt như:
\begin{itemize}
    \item Phage có khả năng chuyển đổi giữa các lối sống (lifestyle switching)
    \item Prophage bị đột biến mất chức năng (defective prophage)
    \item Phage với các cơ chế lysogen hóa không điển hình
    \item Tác động của môi trường chủ lên biểu hiện lối sống
\end{itemize}

Các giới hạn và giả định này được thiết lập để đảm bảo nghiên cứu tập trung và khả thi trong phạm vi luận văn thạc sĩ, đồng thời cung cấp nền tảng vững chắc cho các nghiên cứu mở rộng trong tương lai.
